\section{Agroclimatology \& Field Hydrology}
\label{sec:hydrology}

All water-balance quantities are computed on a \emph{monthly} time step over the months each crop actually occupies, using a per-parcel crop calendar (Section~\ref{sec:calendar}). This replaces the fixed warm/cool half-year windows of Version~1.0, which assigned autumn-sown crops (winter wheat, barley, rapeseed) and perennials (alfalfa) to a single half-year and therefore omitted their spring peak water use.

\subsection{Monthly Reference Evapotranspiration \texorpdfstring{$ET_0$}{ET0}}
\label{sec:et0}
\paragraph{Primary method: monthly FAO-56 Penman--Monteith.}
NASA POWER 1991--2020 monthly climatology provides every variable required by the FAO-56 Penman--Monteith equation~\cite{allen1998fao56,nasapower2024}: mean daily maximum and minimum air temperature at 2~m (\texttt{T2M\_MAX}, \texttt{T2M\_MIN}), relative humidity (\texttt{RH2M}), dew-point temperature (\texttt{T2MDEW}), wind speed at 2~m (\texttt{WS2M}) and all-sky shortwave radiation (\texttt{ALLSKY\_SFC\_SW\_DWN}). Monthly $ET_{0,m}$ ($\mathrm{mm\,day^{-1}}$) is therefore computed with FAO-56 Eq.~6:
\begin{equation}
  ET_0 = \frac{0.408\,\Delta\,(R_n - G) + \gamma\,\dfrac{900}{T_{\mathrm{mean}} + 273}\, u_2 \,(e_s - e_a)}{\Delta + \gamma\,(1 + 0.34\, u_2)},
\end{equation}
with atmospheric pressure $P_{\mathrm{atm}} = 101.3 \bigl(\frac{293 - 0.0065 z}{293}\bigr)^{5.26}$~($\mathrm{kPa}$, Eq.~7), psychrometric constant $\gamma = 0.665\times 10^{-3} P_{\mathrm{atm}}$ (Eq.~8), $e_s = \tfrac{1}{2}[e^{\circ}(T_{\max}) + e^{\circ}(T_{\min})]$ (Eq.~12), and actual vapour pressure from dew point, $e_a = e^{\circ}(T_{\mathrm{dew}})$ (Eq.~14), falling back to $e_a = e_s\,\mathrm{RH}_{\mathrm{mean}}/100$ (Eq.~19) only when dew point is unavailable. Net radiation is $R_n = (1-\alpha)R_s - R_{nl}$ with $\alpha = 0.23$, and net longwave radiation uses the \emph{actual} relative shortwave radiation (FAO-56 Eq.~39):
\begin{equation}
  R_{nl} = \sigma \left[\frac{T_{\max,K}^4 + T_{\min,K}^4}{2}\right]\bigl(0.34 - 0.14\sqrt{e_a}\bigr)\left(1.35\,\frac{R_s}{R_{so}} - 0.35\right), \qquad R_{so} = (0.75 + 2\times10^{-5} z)\,R_a,
\end{equation}
with $R_s/R_{so}$ limited to $\le 1.0$. Monthly soil heat flux is not neglected: $G_m = 0.07\,(T_{\mathrm{mean},m+1} - T_{\mathrm{mean},m-1})$ (FAO-56 Eq.~43; climatological months are cyclic, so the following month is always known). Extraterrestrial radiation $R_a$ on the mid-month day of year $J_m$ follows FAO-56 Eqs.~21--25:
\begin{align}
  d_r &= 1 + 0.033 \cos\!\left(\frac{2\pi}{365} J_m\right), &
  \delta &= 0.409 \sin\!\left(\frac{2\pi}{365} J_m - 1.39\right), \\
  \omega_s &= \arccos\!\bigl(\operatorname{clamp}(-\tan\phi \tan\delta,\, -1,\, 1)\bigr), &
  R_a &= \frac{24(60)}{\pi} G_{sc}\, d_r \Bigl(\omega_s \sin\phi\sin\delta + \cos\phi\cos\delta\sin\omega_s\Bigr),
\end{align}
where $G_{sc} = 0.0820~\mathrm{MJ\,m^{-2}\,min^{-1}}$ and $\phi$ is latitude (rad). The same kernel is used on a daily step (\texttt{calculate\_\allowbreak fao56\_\allowbreak penman\_\allowbreak monteith\_et0}) when daily series are supplied, with $G \approx 0$ (FAO-56 Eq.~42).

\paragraph{Fallback: Hargreaves--Samani.}
Only when humidity, wind or radiation are missing is $ET_{0,m}$ estimated with FAO-56 Eq.~52~\cite{hargreaves1985reference,allen1998fao56}:
\begin{equation}
  ET_{0,m} = 0.0023 \, \bigl(T_{\mathrm{mean},m} + 17.8\bigr) \sqrt{\max\!\bigl(0.1,\, T_{\max,m} - T_{\min,m}\bigr)} \, \bigl(0.408\, R_a\bigr).
\end{equation}
The provenance metadata records which method was used. FAO-56 notes that Eq.~52 tends to over-predict under high relative humidity and under-predict under high wind; both conditions occur at the demonstration sites (humid Po plain near Turin; windy Bosphorus coast near Istanbul), which is why Penman--Monteith is the primary method.

\subsection{Crop Calendar and Monthly Crop Coefficients}
\label{sec:calendar}
\paragraph{Calendar.} Sowing and harvest months for each crop are resolved with the same three-tier precedence as other parameters: farmer-reported dates (Tier~1), otherwise the GGCMI Phase~3 global crop calendar at $0.5^{\circ}$ (planting and maturity day per crop, rainfed and irrigated separately)~\cite{jagermeyr2021calendar}, otherwise catalog defaults. A crop occupies month $m$ with fraction $f_m \in [0,1]$ (partial first and last months).

\paragraph{Single crop coefficient curve.} Crop evapotranspiration uses the FAO-56 \emph{single} crop coefficient $K_c$ (Table~12), not the basal coefficient $K_{cb}$, because the lumped monthly balance does not model soil evaporation separately. Stage lengths $L_{\mathrm{ini}}, L_{\mathrm{dev}}, L_{\mathrm{mid}}, L_{\mathrm{late}}$ (days) come from FAO-56 Table~11 scaled to the resolved season length. Following FAO-56, $K_c$ is constant at $K_{c,\mathrm{ini}}$ during the initial stage, rises linearly to $K_{c,\mathrm{mid}}$ during development, is constant during mid-season and falls linearly to $K_{c,\mathrm{end}}$ during late season. The monthly coefficient $K_{c,m}$ is the mean of this daily curve over the crop days in month~$m$. $K_{c,\mathrm{mid}}$ is adjusted for climate with FAO-56 Eq.~62,
\begin{equation}
  K_{c,\mathrm{mid}} = K_{c,\mathrm{mid}}^{\mathrm{Tab}} + \bigl[0.04\,(u_2 - 2) - 0.004\,(\mathrm{RH}_{\min} - 45)\bigr]\left(\frac{h}{3}\right)^{0.3},
\end{equation}
where $h$ is mean crop height (m) and $\mathrm{RH}_{\min} = 100\, e^{\circ}(T_{\mathrm{dew}})/e^{\circ}(T_{\max})$ during mid-season. Perennial forage (alfalfa) uses the FAO-56 averaged cutting-cycle coefficient for every month of the stand. Monthly crop evapotranspiration is
\begin{equation}
  ET_{c,m} = K_{c,m}\, ET_{0,m}\, N_m\, f_m ,
\end{equation}
with $N_m$ the number of days in month $m$.

\subsection{Effective Rainfall}
Monthly effective rainfall uses the USDA-SCS formula as implemented in FAO CROPWAT~\cite{smith1992cropwat}, applied to each \emph{monthly} total $P_m$ (mm):
\begin{equation}
  P_{\mathrm{eff},m} =
  \begin{cases}
    P_m\,\dfrac{125 - 0.2\,P_m}{125}, & P_m \le 250~\mathrm{mm}, \\[6pt]
    125 + 0.1\,P_m, & P_m > 250~\mathrm{mm}.
  \end{cases}
\end{equation}
Version~1.0 used $\max(0, 0.75\,P_{\mathrm{seas}} - 10)$ on half-year totals. That expression is not the USDA-SCS method, and its offset is meaningful only for monthly totals.

\subsection{Monthly Root-Zone Water Balance}
Total and readily available water follow FAO-56 Eqs.~82--83:
\begin{equation}
  \mathrm{TAW} = 1000\,(\theta_{\mathrm{fc}} - \theta_{\mathrm{wp}})\, Z_r, \qquad
  \mathrm{RAW} = p\,\mathrm{TAW}, \qquad p = p_{\mathrm{Tab}} + 0.04\,(5 - ET_c),
\end{equation}
with rooting depth $Z_r$ and depletion fraction $p_{\mathrm{Tab}}$ from FAO-56 Table~22 ($p$ bounded to $[0.1, 0.8]$; $ET_c$ in $\mathrm{mm\,day^{-1}}$). Field capacity and wilting point are estimated from local or SoilGrids sand, clay and organic matter with the Saxton--Rawls pedotransfer functions~\cite{saxton2006soil}, falling back to FAO-56 Table~19 texture-class midpoints.

Root-zone storage $S$ (mm, $0 \le S \le \mathrm{TAW}$) is carried month by month and across crop years. With start-of-month depletion $D_{r,m} = \mathrm{TAW} - S_{m-1}$, the FAO-56 water-stress coefficient (Eq.~84) is applied to the \emph{root-zone depletion state}, as FAO-56 defines it:
\begin{equation}
  K_{s,m} =
  \begin{cases}
    1, & D_{r,m} \le \mathrm{RAW}, \\[4pt]
    \dfrac{\mathrm{TAW} - D_{r,m}}{(1 - p)\,\mathrm{TAW}}, & D_{r,m} > \mathrm{RAW},
  \end{cases}
  \qquad
  ET_{a,m} = \min\!\bigl(K_{s,m}\, ET_{c,m},\; S_{m-1} + P_{\mathrm{eff},m} + I_{\mathrm{net},m}\bigr),
\end{equation}
\begin{equation}
  S_m = \min\!\bigl(\mathrm{TAW},\; S_{m-1} + P_{\mathrm{eff},m} + I_{\mathrm{net},m} - ET_{a,m}\bigr),
\end{equation}
water above $\mathrm{TAW}$ being treated as drainage. In months without a crop, the parcel is treated as bare or residue-covered soil with $K_{c} = K_{c,\mathrm{ini}}$ of the following crop, which represents mainly soil evaporation in FAO-56.

\paragraph{Initial state.} Year~1 starts from the state obtained by running the balance through the 12 months preceding sowing with the farmer-reported Year~0 crop. Without crop history, $S_0 = 1000\,\bigl(\min(\theta_0, \theta_{\mathrm{fc}}) - \theta_{\mathrm{wp}}\bigr)^+ Z_r$, with $\theta_0$ the latest probe reading taken within 30 days of sowing, or else the NASA SMAP Level-4 root-zone soil moisture climatology for the month before sowing. The ad-hoc factor $0.4$ of Version~1.0 is removed.

\paragraph{Irrigation.} When irrigation is enabled, net irrigation in month $m$ refills the root zone to the readily-available threshold, $I_{\mathrm{net},m} = \max\bigl(0,\, ET_{c,m} - P_{\mathrm{eff},m} - (S_{m-1} - (\mathrm{TAW} - \mathrm{RAW}))\bigr)$, limited by the remaining annual quota. Gross withdrawal is $I_{\mathrm{gross},m} = I_{\mathrm{net},m}/e_a$. The field application efficiency $e_a$ depends on the irrigation method: $0.60$ for surface (border, furrow, basin), $0.75$ for sprinkler and $0.90$ for drip~\cite{fao_irrigation_efficiency}. The annual cap is $\sum_m I_{\mathrm{gross},m} \le Q_{\mathrm{mm}} = Q_{\mathrm{yr}}/(10\,A_{\mathrm{ha}})$, shared by every crop grown in that calendar year. Gravity-fed surface irrigation from canals is common in the Turin area, so declaring the system type matters: at $e_a = 0.60$ the same quota delivers about 30\% less water to the root zone than the single value $0.85$ used in Version~1.0.

\subsection{Seasonal Yield Response to Water (FAO-33)}
Water-limited yield follows the FAO-33 seasonal relationship between relative yield and relative evapotranspiration~\cite{doorenbos1979fao33}, evaluated with the seasonal sums of the monthly balance:
\begin{equation}
  1 - \frac{Y_a}{Y_m} = K_y \left(1 - \frac{\sum_m ET_{a,m}}{\sum_m ET_{c,m}}\right), \qquad
  \mu_{\mathrm{water}} = \max\!\left(0,\; 1 - K_y \left(1 - \frac{\sum_m ET_{a,m}}{\sum_m ET_{c,m}}\right)\right),
\end{equation}
with crop-specific seasonal $K_y$ from FAO-33. The linear relationship was derived for moderate deficits; when the relative evapotranspiration deficit exceeds $0.5$ the result is flagged as \emph{outside the calibrated range} instead of being silently clamped. Version~1.0 substituted the seasonal unmet demand for the daily depletion state in $K_s$ and then applied $K_y(1-K_s)$, which roughly doubled the predicted loss for moderate deficits; the arbitrary floor ($K_s \ge 0.30$) and cap ($\Delta Y \le 0.65$) that compensated for this are removed.

Observed daily climate extremes (\texttt{frost\_\allowbreak days\_\allowbreak count}, \texttt{heatwave\_\allowbreak days\_\allowbreak count}, \texttt{longest\_\allowbreak dry\_\allowbreak spell\_\allowbreak days}) and ESP32 canopy temperature/humidity further modulate cold-kill pruning and heat-stress yield penalties (Section~\ref{sec:optimization}). A seasonal balance cannot resolve a short drought at flowering; this limitation is stated in Section~\ref{sec:limitations}.
